\chapter{정리와 다음 단계}
\label{ch:wrapup}

축하한다. 여기까지 따라왔다면 독자는 이제 토크나이저부터 텍스트 생성까지 LLM의 모든 구성 요소를 직접 구현해 본 것이다. 이 장에서는 1권을 정리하고, 2권 \textit{Make LLM-advanced}에서 다룰 고급 주제를 소개한다.

\section{1권 요약}

\subsection{구현한 것}

\begin{center}
\begin{tabularx}{\textwidth}{lX}
\toprule
\textbf{장} & \textbf{구현 내용} \\
\midrule
3장 & BPE, Unigram 토크나이저 (학습/인코딩/디코딩) \\
4장 & Token/Positional Embedding, EmbeddingStack \\
5장 & Scaled Dot-Product Attention, Causal Mask, Multi-Head Attention \\
6장 & FeedForward (GELU), Transformer Block (Pre-LN) \\
7장 & GPT 모델, LMDataset, Cross-Entropy, AdamW, Cosine LR, Trainer \\
8장 & Greedy/Temperature/Top-K/Top-P 샘플러, generate 함수 \\
\bottomrule
\end{tabularx}
\end{center}

\subsection{테스트}

총 48개의 단위 테스트가 모두 통과한다. 형상 테스트, 수치 테스트, 통합 테스트를 통해 모든 모듈이 올바르게 동작함을 검증했다.

\begin{lstlisting}[style=bashstyle]
$ pytest tests/ -v
========================= 48 passed in 2.47s =========================
\end{lstlisting}

\subsection{독자가 가진 것}

이제 독자는 다음을 할 수 있다:
\begin{itemize}
  \item 자신만의 말뭉치로 토크나이저 학습
  \item 임의의 크기 GPT 모델 생성 (config만 바꾸면 됨)
  \item CPU/GPU에서 학습 루프 실행
  \item 다양한 샘플링 전략으로 텍스트 생성
  \item HuggingFace Transformers의 코드를 읽을 때 내부 동작 이해
\end{itemize}

\section{1권의 한계}

솔직히 말하면, 1권에서 만든 모델은 ChatGPT와는 거리가 멀다. 한계는 크게 세 가지다.

\textbf{크기}: 1권 모델은 1천만$\sim$1억 파라미터 수준이다. GPT-3는 1,750억, GPT-4는 추정 1.7조 파라미터다. 이 크기 차이를 극복하려면 분산 학습이 필수다.

\textbf{아키텍처}: 1권은 2017년 원본 트랜스포머에 가깝다. 현대 LLM(LLaMA, Mistral, GPT-4)은 RoPE, GQA, SwiGLU, RMSNorm 등 더 효율적이고 안정적인 컴포넌트를 사용한다.

\textbf{정렬}: 1권 모델은 그냥 ``다음 토큰 예측''만 학습했다. 사용자의 질문에 답하는 챗봇으로 만들려면 SFT(Supervised Fine-Tuning), RLHF(Reinforcement Learning from Human Feedback), DPO(Direct Preference Optimization) 같은 정렬 기법이 필요하다.

\section{2권에서 다룰 것}

2권 \textit{Make LLM-advanced}는 위 세 가지 한계를 극복하는 법을 다룬다.

\subsection{분산 학습 (2$\sim$4장)}

\begin{itemize}
  \item \textbf{DDP (Distributed Data Parallel)}: 여러 GPU에서 같은 모델 복제 학습.
  \item \textbf{FSDP (Fully Sharded Data Parallel)}: 모델을 샤드로 분할하여 메모리 절약.
  \item \textbf{ZeRO-1/2/3}: 옵티마이저/그래디언트/파라미터를 샤딩.
  \item \textbf{Pipeline Parallel}: 모델을 층별로 다른 GPU에 배치.
  \item \textbf{Tensor Parallel}: 단일 레이어를 여러 GPU로 분할 (Megatron-LM).
\end{itemize}

\subsection{최신 아키텍처 (5$\sim$6장)}

\begin{itemize}
  \item \textbf{RoPE (Rotary Position Embedding)}: 복소수 회전으로 위치 인코딩. extrapolation 우수.
  \item \textbf{GQA (Grouped Query Attention)}: KV 헤드를 줄여 추론 메모리 절약.
  \item \textbf{SwiGLU}: GLU 변종으로 FFN 성능 향상. LLaMA가 사용.
  \item \textbf{RMSNorm}: LayerNorm의 경량화 버전.
  \item \textbf{MoE (Mixture-of-Experts)}: 전문가 네트워크 풀에서 토큰별 상위 K개만 활성화.
\end{itemize}

\subsection{양자화와 추론 최적화 (7장)}

\begin{itemize}
  \item \textbf{INT8/INT4 양자화}: 가중치를 저비트 정수로 변환하여 메모리 4$\sim$8배 절약.
  \item \textbf{AWQ}: 활성값 통계를 활용한 정확도 손실 최소화.
  \item \textbf{GPTQ}: 2차 정보를 사용한 양자화.
  \item \textbf{KV Cache}: 이전 토큰의 K, V를 재사용하여 추론 가속.
  \item \textbf{PagedAttention}: vLLM의 페이지 단위 KV 관리로 메모리 단편화 해결.
\end{itemize}

\subsection{정렬과 파인튜닝 (8장)}

\begin{itemize}
  \item \textbf{SFT (Supervised Fine-Tuning)}: 지도학습으로 챗봇 행동 학습.
  \item \textbf{RLHF}: 보상 모델 + PPO로 인간 선호도 최적화.
  \item \textbf{DPO}: RLHF의 복잡함 없이 직접 선호도 최적화.
  \item \textbf{LoRA}: 저랭크 어댑터로 효율적 파인튜닝.
\end{itemize}

\subsection{데이터 파이프라인 (9장)}

\begin{itemize}
  \item \textbf{품질 필터}: 길이, 언어, 반복 비율로 필터링.
  \item \textbf{MinHash 중복 제거}: Jaccard 유사도 근사로 중복 문서 탐지.
  \item \textbf{합성 데이터}: 템플릿 기반 학습 데이터 생성.
\end{itemize}

\section{추천 학습 경로}

2권으로 넘어가기 전에 추천하는 활동:

\begin{enumerate}
  \item \textbf{작은 프로젝트}: 1권 코드로 자신의 관심 분야 텍스트(예: 위키백과 한국어, GitHub 코드)를 학습해보라. 10만 토큰 정도면 1시간 안에 학습된다.
  \item \textbf{논문 읽기}: ``Attention Is All You Need''(Vaswani et al. 2017), ``Language Models are Few-Shot Learners''(Brown et al. 2020)를 읽어보라. 1권을 읽었으므로 훨씬 이해가 잘 될 것이다.
  \item \textbf{HuggingFace 코드 읽기}: \code{transformers/models/gpt2/modeling\_gpt2.py}를 열어보라. 우리가 만든 코드와 얼마나 비슷한지, 어떤 차이가 있는지 비교하라.
\end{enumerate}

\section{마지막으로}

LLM은 복잡해 보이지만, 결국 토크나이저, 임베딩, 어텐션, FFN, 학습 루프의 조합이다. 이 기본기를 익히면 어떤 새로운 아키텍처, 어떤 새로운 최적화 기법이 나와도 이해할 수 있다. 2권에서는 이 기반 위에 ChatGPT 수준의 기술을 쌓아올릴 것이다.

코딩을 멈추지 마라. 실험을 멈추지 마라. 그것이 실력을 키우는 유일한 길이다.

\begin{flushright}
\textit{1권 끝}\\[0.5em]
\textbf{dirmich}
\end{flushright}
